\begin{lemma}
\label{lemma-structure-torsion-Frobenius-regular}
\begin{reference}
Follows from \cite[Corollary 3.6]{Huneke-Sharp} with a
little bit of work. Also follows directly from
\cite[Theorem 1.4]{Lyubeznik2}.
\end{reference}
Let $p$ be a prime number. Let $(A, \mathfrak m, k)$
be a regular local ring with $p = 0$. Denote $F : A \to A$, $a \mapsto a^p$
be the Frobenius endomorphism. Let $M$ be a $\mathfrak m$-power torsion module
such that $M \otimes_{A, F} A \cong M$. Then $M$ is isomorphic to a direct sum
of copies of the injective hull $E$ of $k$.
\end{lemma}

\begin{proof}
Choose a set $J$ and an $A$-module homorphism
$\varphi : M \to \bigoplus_{j \in J} E$ which maps
$M[\mathfrak m]$ isomorphically onto
$(\bigoplus_{j \in J} E)[\mathfrak m] = \bigoplus_{j \in J} k$.
We claim that $\varphi$ is an isomorphism, i.e., bijective.

\medskip\noindent
Injective. Let $z \in M$ be nonzero. Since $M$ is $\mathfrak m$-power torsion
we can choose an element $f \in A$ such that $fz \in M[\mathfrak m]$ and
$fz \not = 0$. Then $\varphi(fz) = f\varphi(z)$ is nonzero, hence
$\varphi(z)$ is nonzero.

\medskip\noindent
Surjective. Recall that $F$ is flat, see
Lemma \ref{lemma-frobenius-flat-regular}.
Let $x_1, \ldots, x_d$ be a minimal system of generators of
$\mathfrak m$. Denote
$$
M_n = M[x_1^{p^n}, \ldots, x_d^{p^n}]
$$
the submodule of $M$ consisting of elements killed by
$x_1^{p^n}, \ldots, x_d^{p^n}$. So $M_0 = M[\mathfrak m]$
is a vector space over $k$. Also $M = \bigcup M_n$ by our
assumption that $M$ is $\mathfrak m$-power torsion. Since $F^n$ is flat and
$F^n(x_i) = x_i^{p^n}$ we have
$$
M_n \cong (M \otimes_{A, F^n} A)[x_1^{p^n}, \ldots, x_d^{p^n}] =
M[x_1, \ldots, x_d] \otimes_{A, F^n} A =
M_0 \otimes_k A/(x_1^{p^n}, \ldots, x_d^{p^n})
$$
Thus $M_n$ is free over $A/(x_1^{p^n}, \ldots, x_d^{p^n})$.
A computation shows that every element of $A/(x_1^{p^n}, \ldots, x_d^{p^n})$
annihilated by $x_1^{p^n - 1}$ is divisible by $x_1$; for example
you can use that $A/(x_1^{p^n}, \ldots, x_d^{p^n}) \cong
k[x_1, \ldots, x_d]/(x_1^{p^n}, \ldots, x_d^{p^n})$ by Algebra, Lemma
\ref{algebra-lemma-regular-complete-containing-coefficient-field}.
Thus the same is true for every element of $M_n$.
Since every element of $M$ is in $M_n$ for all $n \gg 0$
and since every element of $M$ is killed by some power of
$x_1$, we conclude that $M$ is $x_1$-divisible.

\medskip\noindent
Let $x = x_1$. Above we have seen that $M$ is $x$-divisible.
If $\varphi$ is not surjective, then we can choose
$e \in \bigoplus_{j \in J} E$ not in $M$.
Arguing as above we may assume $\mathfrak m e \subset M$,
in particular $x e \in M$. There exists an element
$z_1 \in M$ with $x z_1 = x e$. Hence
$x(z_1 - e) = 0$. Replacing $e$ by $e - z_1$
we may assume $e$ is annihilated by $x$.
Thus it suffices to prove that
$$
\varphi[x] :
M[x]
\longrightarrow
\left(\bigoplus\nolimits_{j \in J} E\right)[x] =
\bigoplus\nolimits_{j \in J} E[x]
$$
is surjective. If $d = 1$, this is true by construction of $\varphi$.
If $d > 1$, then we observe that $E[x]$ is the injective hull
of the residue field of the regular ring $A/xA$, see
Dualizing Complexes, Lemma \ref{dualizing-lemma-quotient}.
Observe that $M[x]$ as a module over $A/xA$
is $\mathfrak m/(x)$-power torsion and we have
\begin{align*}
M[x] \otimes_{A/xA, F} A/xA
& = M[x] \otimes_{A, F} A \otimes_A A/xA \\
& = (M \otimes_{A, F} A)[x^p] \otimes_A A/xA \\
& \cong M[x^p] \otimes_A A/xA
\end{align*}
Argue using flatness of $F$ as before.
We claim that $M[x^p] \otimes_A A/xA \to M[x]$,
$z \otimes 1 \mapsto x^{p - 1}z$ is an isomorphism.
This can be seen by proving it for
each of the modules $M_n$, $n > 0$ defined above
where it follows by the same result for
$A/(x_1^{p^n}, \ldots, x_d^{p^n})$ and $x = x_1$.
Thus by induction on $\dim(A)$ we conclude that $\varphi[x]$
is surjective as desired.
\end{proof}